\documentclass[11pt,a4paper]{article}
\usepackage{fontspec,amsmath,amssymb,amsthm,microtype}
\setmainfont{lmroman10-regular.otf}[BoldFont=lmroman10-bold.otf,ItalicFont=lmroman10-italic.otf,BoldItalicFont=lmroman10-bolditalic.otf]
\usepackage{xcolor}
\usepackage[most]{tcolorbox}
\usepackage{mathrsfs}
\providecommand{\lt}{<}
\providecommand{\gt}{>}
\setmonofont{DejaVuSansMono.ttf}[BoldFont=DejaVuSansMono-Bold.ttf,ItalicFont=DejaVuSansMono-Oblique.ttf,BoldItalicFont=DejaVuSansMono-BoldOblique.ttf,Scale=0.86]
\newfontfamily\leanfallback{STIXTwoMath-Regular.otf}[Scale=0.86]
\definecolor{leanrule}{gray}{0.62}
\definecolor{leanback}{gray}{0.962}
\definecolor{leantitle}{gray}{0.22}
\definecolor{leanmuted}{gray}{0.38}
\newtcolorbox{leanbox}[1]{enhanced,breakable,colback=leanback,colframe=leanrule,
  boxrule=0.4pt,arc=0.6mm,left=2.4mm,right=2.4mm,top=1.8mm,bottom=1.8mm,
  before skip=0.8em,after skip=1em,fonttitle=\small\bfseries,coltitle=leantitle,
  colbacktitle=leanback,titlerule=0.3pt,title={#1}}
\newlength{\leancharwidth}
\newenvironment{leancode}{\par\smallskip\begingroup\ttfamily\footnotesize\raggedright
  \settowidth{\leancharwidth}{M}\setlength{\parskip}{0pt}\setlength{\parindent}{0pt}
  \frenchspacing}{\par\endgroup\smallskip}
\newcommand{\leanline}[2]{\hangindent=\dimexpr#1\leancharwidth+4\leancharwidth\relax
  \hangafter=1\noindent\hspace*{#1\leancharwidth}#2\par}
\newcommand{\leansource}[1]{\par{\footnotesize\color{leanmuted}#1}\par}
\newcommand{\leanlabel}[1]{\par\smallskip\noindent{\small\bfseries #1}\par}
\usepackage[margin=25mm]{geometry}
\usepackage[unicode,hidelinks]{hyperref}
\hypersetup{pdfauthor={LangevinLean Project}}
\setlength{\emergencystretch}{3em}
\setlength{\parskip}{0.4em}
\setlength{\parindent}{0pt}
\allowdisplaybreaks[2]
\raggedbottom
\title{\texorpdfstring{The semigroup law from the actual conditional Markov identity}{The semigroup law from the actual conditional Markov identity}}
\author{LangevinLean Project}
\date{}
\begin{document}\maketitle
Draft mathematical exposition; these notes may contain mistakes.


The conditional Markov identity on the actual initial/\allowbreak{}Brownian past implies the composition law for the canonical real observable contractions. Its proof uses actual adaptedness of time-zero evaluation and the defining adjoint pairings.


\section{\texorpdfstring{Conditioning the actual two-time pairing}{Conditioning the actual two-time pairing}}


\textbf{Claim.} If \(\mathbb E[U_{a+t}f\mid\mathcal F_a]=U_aP_tf\), then \(P_{a+t}=P_aP_t\).



\textbf{Proof.} Time-zero evaluation is measurable with respect to the actual past because that past contains the initialization. Pair the conditional identity with \(U_0g\). The conditional-expectation projection preserves its pairing with this adapted vector. The adjoint compression identities identify the resulting sides with \(\langle P_{a+t}f,g\rangle\) and \(\langle P_aP_tf,g\rangle\). Equality against all test observables proves the operator identity.

\begin{leanbox}{Lean: \texttt{evaluation\_\allowbreak{}zero\_\allowbreak{}adapted\_\allowbreak{}pastAt}, \texttt{op\_\allowbreak{}add\_\allowbreak{}of\_\allowbreak{}markov}}
\leansource{\texttt{evaluation\_\allowbreak{}zero\_\allowbreak{}adapted\_\allowbreak{}pastAt} \quad WassersteinMixing/\allowbreak{}SemigroupComposition.lean, lines 24–39}
\begin{leancode}
\leanline{0}{theorem\ evaluation\_zero\_adapted\_pastAt\ (P\ :\ Parameters)\ (V\ :\ E\ P.dimension\ →\ ℝ)}
\leanline{4}{(hV\ :\ PotentialAssumptions\ P\ V)\ (a\ :\ ℝ)\ (f\ :\ Lp\ F\ 2\ (ULA.gibbsMeasure\ V))\ :}
\leanline{4}{AEStronglyMeasurable[ULA.lpLangevinPast\ (actualProbability\ P\ V)}
\leanline{6}{(initial\ P\ V\ hV)\ (actualBrownianNoise\ P\ V\ hV)\ a]}
\leanline{6}{(evaluation\ ({\leanfallback 𝕜}\ :=\ ℝ)\ P\ V\ hV\ 0\ f)\ (actualProbability\ P\ V)}
\end{leancode}
\leansource{\texttt{op\_\allowbreak{}add\_\allowbreak{}of\_\allowbreak{}markov} \quad WassersteinMixing/\allowbreak{}SemigroupComposition.lean, lines 44–66}
\begin{leancode}
\leanline{0}{theorem\ op\_add\_of\_markov\ (P\ :\ Parameters)\ (V\ :\ E\ P.dimension\ →\ ℝ)}
\leanline{4}{(hV\ :\ PotentialAssumptions\ P\ V)\ (a\ t\ :\ ℝ)}
\leanline{4}{(hmarkov\ :\ ∀\ f\ :\ Lp\ F\ 2\ (ULA.gibbsMeasure\ V),}
\leanline{6}{(condExpL2\ F\ ℝ\ (ULA.lpLangevinPast\_le\ (actualProbability\ P\ V)}
\leanline{8}{(initial\ P\ V\ hV)\ (actualBrownianNoise\ P\ V\ hV)\ a)}
\leanline{8}{(evaluation\ ({\leanfallback 𝕜}\ :=\ ℝ)\ P\ V\ hV\ (a\ +\ t)\ f)\ :\ Lp\ F\ 2\ (actualProbability\ P\ V))\ =}
\leanline{8}{evaluation\ ({\leanfallback 𝕜}\ :=\ ℝ)\ P\ V\ hV\ a\ (op\ ({\leanfallback 𝕜}\ :=\ ℝ)\ P\ V\ hV\ t\ f))\ :}
\leanline{4}{op\ ({\leanfallback 𝕜}\ :=\ ℝ)\ (F\ :=\ F)\ P\ V\ hV\ (a\ +\ t)\ =\ op\ P\ V\ hV\ a\ ∘L\ op\ P\ V\ hV\ t}
\end{leancode}
\leanlabel{How the Lean proof follows the argument}
\texttt{evaluation\_\allowbreak{}zero\_\allowbreak{}adapted\_\allowbreak{}past\allowbreak{}At} uses the actual initial-state generator of the past. \texttt{op\_\allowbreak{}add\_\allowbreak{}of\_\allowbreak{}markov} applies the real conditional-expectation pairing theorem and the exact evaluation compression identities, then uses inner-product extensionality.
\end{leanbox}


Earlier results used here: \cite[Section~3.26.3]{LangevinLean2026ULASC}; \cite[Section~3.54.4]{LangevinLean2026ULASC}; \cite[Section~3.54.5]{LangevinLean2026ULASC}; \cite[Section~3.54.6]{LangevinLean2026ULASC}.

\begin{thebibliography}{ULA-LSI}
\bibitem[ULA-SC]{LangevinLean2026ULASC}
LangevinLean Project. \emph{Unadjusted Langevin under Strong Convexity}. LL-01, 2026. Draft report, content revision d34ead06987c. \url{https://odakyildiz.github.io/LangevinLean/docs/pdf/UnadjustedLangevin/FullProof.pdf}.
\end{thebibliography}

\end{document}
